\section{Related work}
\label{sec:related}

\paragraph{Decomposed visual decisions.} Attribute-based recognition~\cite{lampert2009attribute}, concept bottleneck models (CBMs)~\cite{koh2020cbm} and program-based VLM pipelines~\cite{gupta2023visprog,suris2023vipergpt} predict through human-readable intermediate facts. \lux shares the decomposition. Our question is different: what verdict supervision can and cannot teach about the facts, and what happens to a trained reader when the composition rule is edited afterwards.

\paragraph{Concept leakage.} CBM concepts can carry information beyond their meaning~\cite{margeloiu2021cbm,mahinpei2021promises}. Havasi \etal~\cite{havasi2022leakage} trace the gap between label and concept accuracy to an insufficient concept set and an inexpressive concept predictor, and close it with a side channel for latent concepts. We measure \emph{verdict leakage}---how much an absent fact's reading depends on the scene's verdict---as a violation of the conditional independence our analysis assumes (\cref{sec:exp-health}), and we repair it without a side channel, which would route rule-specific information back into the decision.

\paragraph{Reasoning shortcuts.} When concepts are learned through a fixed knowledge base, several concept assignments can fit the labels equally well~\cite{marconato2023not}; concept semantics are identifiable only without joint shortcuts~\cite{bortolotti2025shortcuts}. rsbench~\cite{bortolotti2024rsbench} exposes shortcuts by editing traffic rules on synthetic driving scenes. Overexposure is the graded form of this phenomenon under monotone, saturating rules: clipped facts are not swapped but \emph{unconstrained}. We contribute an information measure in place of shortcut counts, a regret bound under rule edits, a repair objective, and real-image evidence.

\paragraph{Discriminative evidence and missing labels.} Weakly supervised localisation activates only the most discriminative region and is repaired by erasing or hiding it~\cite{zhou2016cam,wei2017erasing,singh2017hide}; multiple-instance learning has unreliable instance labels~\cite{ilse2018attention}; gradient starvation leaves predictive features unlearned~\cite{pezeshki2021gradient}. These are overexposure in space, in bags and in features. Incomplete fact annotations relate to learning from positive and unlabelled data~\cite{elkan2008pu} and from single positive labels~\cite{cole2021single}; instead of constraining the expected number of positives, we identify the annotation model on the cells whose presence the verdict pins down.
